\section{Methodology}
\label{sec:methodology}

\subsection{DEM simulation of hollow cylinder tests}
\label{subsec:dem_hca}

DEM simulations were performed in \textsc{PFC3D} \citep{Itasca2021} to reproduce undrained cyclic torsional shear in an HCA. Unlike periodic-cell element tests, the present HCA model uses rigid boundaries (inner and outer cylinders, top and bottom planes) and six torsional blades to apply shear while preserving laboratory-relevant boundary conditions.

The specimen geometry (Fig.~\ref{fig:specimen_generation}) follows a laboratory-scale HCA configuration with inner diameter \SI{6}{cm}, outer diameter \SI{10}{cm}, and specimen height \SI{10}{cm} \citep{Vargas2020}. A rolling-resistance contact model \citep{IwashitaOda1998,WensrichKatterfeld2012} is adopted to capture non-spherical grain effects, with a particle size distribution in the range of \SIrange{1.0}{3.0}{mm}, which is a 10$\times$ scaled-up analogue of the grain size distribution of Toyoura sand, a scaling approach commonly adopted in DEM simulations of this sand to reduce particle count while preserving the non-dimensional grading \citep{Song2024}. The DEM parameters are summarized in Table~\ref{tab:dem_parameters}. Particles are generated by gravity deposition from the upper part of the apparatus. Gravity is then removed to ensure a uniform stress state, and six torsional blades, each 15~mm tall, are inserted into the specimen (Fig.~\ref{fig:specimen_generation}b). Three blades at each end of the specimen extend radially across the annulus between the inner and outer cylinders, uniformly spaced around the specimen axis. Compared with the conventional approach of constraining boundary particles to transmit torque \citep[e.g.,][]{LiZhangGutierrez2015,Liu2021}, these independent wall-based blades avoid interference with the radial expansion or contraction of the cylinders, thereby preserving accurate radial stress measurement. The cylindrical walls and end platens are assigned zero friction ($\mu_{pw}=0$, Table~\ref{tab:dem_parameters}) so that they act purely as servo-controlled pressure boundaries; torque is transmitted geometrically through normal contact forces on the blade faces and does not rely on a wall-friction coefficient.

\begin{table}[t]
  \centering
  \caption{Parameters in DEM simulation.}
  \label{tab:dem_parameters}
  \begin{tabular}{@{}lr@{}}
    \toprule
    Description & Value \\
    \midrule
    Number of particles, $N_p$                          & 53{,}764 \\
    Particle density, $\rho_s$ (\si{kg/m^3})            & 2650 \\
    Young's modulus, $E$ (\si{GPa})                     & 0.3 \\
    Normal-to-shear stiffness ratio, $\kappa$           & 2.0 \\
    Interparticle friction, $\mu_{pp}$                  & 0.0\textsuperscript{a}\,/\,0.50\textsuperscript{b} \\
    Particle--wall friction, $\mu_{pw}$                 & 0.0 \\
    Normal critical damping ratio, $\beta_n$            & 0.7 \\
    Shear critical damping ratio, $\beta_s$             & 0.5 \\
    Rolling friction coefficient, $\mu_r$               & 0.1 \\
    \bottomrule
    \multicolumn{2}{@{}l}{\footnotesize \textsuperscript{a}Isotropic consolidation stage.} \\
    \multicolumn{2}{@{}l}{\footnotesize \textsuperscript{b}Anisotropic consolidation and cyclic shear stages.}
  \end{tabular}
\end{table}

\begin{figure*}
  \centering
  \includegraphics[width=\textwidth]{specimen_generation_standalone.pdf}
  \caption{Initial specimen generation process in the DEM-based hollow cylinder apparatus model.
    Particles are colored by radius.
    (a)~Particle generation modeling air-pluviation.
    (b)~Insertion of torsional blades.}
  \label{fig:specimen_generation}
\end{figure*}

The stress and strain evaluation expressions for the HCA geometry are summarized in Appendix~\ref{sec:appendix_hca}.

\subsection{Specimen preparation and anisotropic consolidation protocol}
\label{subsec:preparation}

Specimens are prepared using an isotropic-consolidation to anisotropic-consolidation (IC--AC) protocol, with the consolidation anisotropy characterized by the lateral-to-vertical effective stress ratio $\Kzero = \sigma_r^{\prime}/\sigma_z^{\prime}$, where $\sigma_r^{\prime}$ and $\sigma_z^{\prime}$ are the radial and vertical effective stresses (Table~\ref{tab:hca_equations}). First, the assembly is compressed under zero interparticle friction to a target void ratio, determined through the iterative calibration procedure described in Section~\ref{subsec:calibration_validation}: approximately 0.736 for the very dense state and 0.742 for the dense state. Because the DEM particles are spherical (in contrast to the subrounded-to-subangular grains of real Toyoura sand) and because no standardised procedure exists for determining the packing limits $(e_{\max}, e_{\min})$ of a DEM assembly, these void ratios are not converted to a relative density through the textbook formula with laboratory Toyoura limits. The labels ``$D_r \approx 90\%$'' and ``$D_r \approx 75\%$'' are instead used as behaviour-matched references to the ``very dense'' ($D_r = 88\!-\!95\%$) and ``dense'' ($D_r = 75\!-\!82\%$) density bands of \citet{Ishihara1985}: the DEM target void ratios are tuned so that the macroscopic monotonic and cyclic responses reproduce those observed in the laboratory at the corresponding $D_r$ conditions (Section~\ref{subsec:calibration_validation}). The interparticle friction coefficient is then restored to 0.5, and the specimen is consolidated from an isotropic state (\Kzero{} = 1.0, $\pmean=\SI{10}{kPa}$) to the target mean effective stress ($\pmean=\SI{100}{kPa}$) along a linear anisotropic stress path.

During the AC stage, $\pmean$ is increased from \SI{10}{kPa} to \SI{100}{kPa} through ten equally spaced target stress states along a linear anisotropic stress path. The resulting stress paths and void-ratio evolution for the three primary cases are shown in Fig.~\ref{fig:ac_path}. The markers in Fig.~\ref{fig:ac_path_a} correspond to the prescribed target states; the solid lines between them represent the actual stress evolution, which exhibits minor fluctuations that arise because different \Kzero{} states require different amounts of axial strain to reach the same $\pmean$ increment---compression-state specimens (\Kzero{} = 0.5) undergo larger axial strain per step, causing the axial stress to lag slightly behind the radial stress during each increment. These fluctuations vanish at each target state, confirming that the servo accurately satisfies the prescribed stress conditions. The final void ratio differs only slightly among the three states (on the order of $10^{-3}$), which allows the comparison to focus on fabric-related effects rather than void-ratio differences. Figure~\ref{fig:ac_specimens} shows the three specimens after consolidation, with force-chain networks and particle arrangements displayed side by side; the visual contrast in force-chain directionality reflects the different \Kzero{} states.

\begin{figure}
  \centering
  \begin{subfigure}[t]{0.432\linewidth}
    \centering
    \safeincludegraphics[width=\linewidth]{anisotropic_consolidation_stress_path.png}
    \caption{$\pmean$--$q$ stress path.}
    \label{fig:ac_path_a}
  \end{subfigure}
  \hfill
  \begin{subfigure}[t]{0.432\linewidth}
    \centering
    \safeincludegraphics[width=\linewidth]{anisotropic_consolidation_void_ratio.png}
    \caption{$\pmean$--$e$ evolution.}
    \label{fig:ac_path_b}
  \end{subfigure}
  \caption{Stress and void-ratio evolution during anisotropic consolidation for the three representative \Kzero{} states.}
  \label{fig:ac_path}
\end{figure}

\begin{figure*}
  \centering
  \begin{tikzpicture}
    % --- (a) K0=0.5 ---
    \node[anchor=north west, inner sep=0] (imgA) at (0,0)
      {\includegraphics[height=4.8cm]{specimen_k0_05_after_ac.png}};
    \node[font=\small, anchor=north] at (imgA.south) {(a) \Kzero{} = 0.5};
    % --- (b) K0=1.0 ---
    \node[anchor=north west, inner sep=0] (imgB) at ([xshift=3pt]imgA.north east)
      {\includegraphics[height=4.8cm]{specimen_k0_10_after_ac.png}};
    \node[font=\small, anchor=north] at (imgB.south) {(b) \Kzero{} = 1.0};
    % --- (c) K0=2.0 ---
    \node[anchor=north west, inner sep=0] (imgC) at ([xshift=3pt]imgB.north east)
      {\includegraphics[height=4.8cm]{specimen_k0_20_after_ac.png}};
    \node[font=\small, anchor=north] at (imgC.south) {(c) \Kzero{} = 2.0};
    % --- Colorbar ---
    \node[anchor=north west, inner sep=0] (cbar) at ([xshift=3pt]imgC.north east)
      {\includegraphics[height=4.8cm]{colorbar_viridis_force.png}};
    % Tick labels
    \foreach \val/\frac in {2.0/0, 1.5/0.25, 1.0/0.5, 0.5/0.75, 0/1.0} {
      \node[font=\footnotesize, anchor=west] at
        ([xshift=2pt, yshift={-\frac*4.8cm}]cbar.north east) {\val};
    }
    % Title
    \node[font=\footnotesize, rotate=90, anchor=south] at
      ([xshift=34pt]cbar.east) {Contact force magnitude (N)};
  \end{tikzpicture}
  \caption{DEM HCA specimens after anisotropic consolidation to $\pmean=\SI{100}{kPa}$.
    The left half of each panel shows interparticle force chains colored by contact force magnitude.}
  \label{fig:ac_specimens}
\end{figure*}

The cyclic shear simulation program consists of two sets. Throughout the parametric \Kzero{} study reported here, the outer and inner cell pressures are held at the post-consolidation value ($p_o=p_i=\sigma_{r0}$, so $p_{dif,r}^{tar}=0$ and the mean radial total stress $\sigma_r$ is time-invariant), and the axial boundary is imposed in displacement-control mode ($dH/dt=0$) matching the protocol of \citet{Ishihara1985}; the constant-total-axial-stress formulation of \cref{subsec:cyclic_loading} is additionally exercised in the verification case of \cref{subsec:numerical_verification}, and the general time-dependent $p_{dif,r}^{tar}(t)$ option remains available in the servo code for future studies that use differential cell control. The cyclic stress ratio is defined as $\mathrm{CSR} = \tau_{cyc}/p_0^{\prime}$, with $\tau_{cyc}$ the amplitude of the sinusoidal shear stress $\tau_{z\theta}(t) = \tau_{cyc}\sin(2\pi f t)$ at loading frequency $f = \SI{8}{Hz}$ (Table~\ref{tab:hca_equations}) and $p_0^{\prime}$ the initial mean effective stress after consolidation. The corresponding inertial number stays well below the strict quasi-static threshold throughout the pre-liquefaction stage; this is verified quantitatively in \cref{subsec:numerical_verification}. In the primary set, three representative \Kzero{} states---0.5 (compression), 1.0 (isotropic), and 2.0 (extension)---are examined at a very dense state ($D_r \approx 90\%$) across multiple CSR levels to characterize how liquefaction resistance varies with the direction and magnitude of stress anisotropy, and to trace the associated energy dissipation, stiffness degradation, and fabric evolution. In the expanded set, a dense state ($D_r \approx 75\%$) is introduced alongside the very dense state, and two additional \Kzero{} values (0.67 and 1.5) are included to refine the coverage of anisotropy levels, yielding ten cases ($2 \times 5$) at a single CSR = 0.300. This expanded set allows examination of whether the \Kzero{}-dependent liquefaction resistance trend observed at $D_r \approx 90\%$ persists at a lower density, and provides a broader parametric basis for the subsequent microscopic analysis. Table~\ref{tab:case_matrix} summarizes the full simulation program and the resulting $N_L$ for each case.

\begin{table}[t]
  \centering
  \caption{Cyclic shear simulation program: CSR levels tested at each combination of relative density and consolidation stress ratio. All cases use the displacement-control axial boundary ($dH/dt=0$) matching \citet{Ishihara1985}, except as noted.}
  \label{tab:case_matrix}
  \begin{tabular}{@{}ccl@{}}
    \toprule
    $D_r$ (\%) & $\Kzero$ & CSR levels \\
    \midrule
    90 & 0.50 & 0.200, 0.250, 0.300, 0.350, 0.400 \\
    90 & 0.67 & 0.300 \\
    90 & 1.00 & 0.200, 0.250, 0.300, 0.350, 0.400 \\
    90 & 1.50 & 0.300 \\
    90 & 2.00 & 0.200, 0.250, 0.300, 0.350, 0.400 \\
    \midrule
    75 & 0.50 & 0.300 \\
    75 & 0.67 & 0.300 \\
    75 & 1.00 & 0.300, 0.400$^{\dagger}$ \\
    75 & 1.50 & 0.300 \\
    75 & 2.00 & 0.300 \\
    \bottomrule
  \end{tabular}\\[2pt]
  \begin{minipage}{0.95\linewidth}
    \footnotesize $^{\dagger}$\,Numerical-verification case (\cref{subsec:numerical_verification}), run under the constant-total-axial-stress axial boundary so that the full $2{\times}2$ coupled servo is exercised.
  \end{minipage}
\end{table}

\subsection{Undrained cyclic torsional shear loading}
\label{subsec:cyclic_loading}

\begin{figure*}
  \centering
  \includegraphics[width=\textwidth]{fig_hca_standalone.pdf}
  \caption{Schematic of the hollow cylinder apparatus (HCA) specimen.
    The left half of each panel shows the laboratory boundary conditions with total stresses
    ($p_o$, $p_i$, $p_z$);
    the right half shows the DEM representation with effective stresses
    ($p_o^\prime$, $p_i^\prime$, $p_z^\prime$).
    (a)~Perspective view with specimen height $H$.
    (b)~Plan view showing the annular geometry with inner radius $r$,
    outer radius $R$, and torque $T$.}
  \label{fig:hca_schematic}
\end{figure*}

In a laboratory HCA (left half of Fig.~\ref{fig:hca_schematic}), the specimen is enclosed between an outer chamber (pressure $p_o$) and an inner chamber (pressure $p_i$), which are controlled through independent channels and can therefore be prescribed as general functions of time. An axial load is applied through the loading piston, producing a vertical pressure $p_z$ that is likewise independently prescribed, and a torque $T$ is transmitted through end platens. Under undrained conditions the pore water pressure $u$, monitored by pressure transducers, is uniform throughout the specimen, so the effective pressures are $p_o^{\prime}=p_o-u$, $p_i^{\prime}=p_i-u$, and $p_z^{\prime}=p_z-u$. Because $u$ drops out in any difference of effective pressures, the radial pressure difference equals its total-stress counterpart, $p_o^{\prime}-p_i^{\prime}=p_o-p_i$, and so does the axial pressure difference between the vertical and average radial effective stresses (Table~\ref{tab:hca_equations}): $p_z^{\prime}-\sigma_r^{\prime}=p_z-\sigma_r$.

The constant-volume method is a well-established approach in DEM for reproducing undrained cyclic loading \citep{YimsiriSoga2010}. In this method, the pore water pressure is not modelled explicitly; instead, the contact forces between particles and walls yield effective stresses ($p_o^{\prime}$, $p_i^{\prime}$, $p_z^{\prime}$; right half of Fig.~\ref{fig:hca_schematic}) directly (Table~\ref{tab:hca_equations}). The DEM-based HCA specimen has four kinematic degrees of freedom---inner radius rate ($dr/dt$), outer radius rate ($dR/dt$), height rate ($dH/dt$), and rotation angle rate ($d\theta/dt$), where $r$ is the inner radius, $R$ the outer radius (Fig.~\ref{fig:hca_schematic}b), and $H$ the height of the hollow cylindrical specimen (Fig.~\ref{fig:hca_schematic}a)---and the following laboratory boundary conditions to reproduce: the undrained (constant-volume) constraint, the inner and outer cell pressures, the axial stress, and the torsional shear stress.

Among these, the constant-volume constraint provides one equation linking $dr/dt$, $dR/dt$, and $dH/dt$, while the torsional condition is satisfied by adjusting the rotation rate $d\theta/dt$ to match the target torque. The remaining challenge is to map the three effective pressures accessible in DEM---$p_i^{\prime}$, $p_o^{\prime}$, and $p_z^{\prime}$---to the laboratory total-stress boundary conditions, where the unknown pore water pressure $u$ prevents a direct one-to-one correspondence. The difference relations established above resolve this difficulty: because $u$ drops out, the two differences $p_o^{\prime}-p_i^{\prime}=p_o-p_i$ and $p_z^{\prime}-\sigma_r^{\prime}=p_z-\sigma_r$ map the DEM effective stresses to the laboratory total-stress conditions without requiring $u$, yielding two independent equations that, together with the volume constraint and the torsional condition, close the system of four equations for four unknowns.

\paragraph{Radial condition.}
The radial condition controls the pressure difference between the outer and inner cylinders. The radial pressure difference and its servo error are defined as:
\begin{equation}
p_{dif,r}^{\prime} = p_o^{\prime} - p_i^{\prime}, \qquad e_r = p_{dif,r}^{\prime} - p_{dif,r}^{tar},
\label{eq:cond1_radial}
\end{equation}
where $p_o^{\prime}$ and $p_i^{\prime}$ are the outer and inner radial effective pressures (Table~\ref{tab:hca_equations}) and $p_{dif,r}^{tar}$ is the target radial pressure difference. The servo must drive $e_r$ to zero at each timestep to satisfy this condition.

\paragraph{Axial condition.}
Laboratory HCA tests can operate under either constant total axial stress or constant specimen height.
(i)~\emph{Stress-control mode}: the servo controls the difference between the vertical and radial effective stresses, analogous to the radial condition in \cref{eq:cond1_radial}. The axial pressure difference and its servo error are defined as:
\begin{equation}
p_{dif,z}^{\prime} = p_z^{\prime} - \sigma_r^{\prime}, \qquad e_z = p_{dif,z}^{\prime} - p_{dif,z}^{tar},
\label{eq:cond2_axial}
\end{equation}
where $p_z^{\prime}$ is the vertical effective pressure (Table~\ref{tab:hca_equations}), $\sigma_r^{\prime}$ is the average radial effective stress, and $p_{dif,z}^{tar}$ is the target axial pressure difference. The servo must drive $e_z$ to zero at each timestep to satisfy this condition.
(ii)~\emph{Displacement-control mode}: the specimen height is held constant ($dH/dt=0$), and $p_{dif,z}^{\prime}$ evolves freely.

\paragraph{Torsional condition.}
The torque error is defined as:
\begin{equation}
e_T = T - T^{tar},
\label{eq:cond3_torsion_error}
\end{equation}
where $T$ is the current torque and $T^{tar}$ is the target torque. The servo must drive $e_T$ to zero at each timestep to satisfy the prescribed torsional loading.

\paragraph{Volume constraint.}
The undrained (constant-volume) condition requires (with the outward-positive sign convention $dr>0$, $dR>0$ denoting radius increases):
\begin{equation}
2\pi H\!\left(R\frac{dR}{dt} - r\frac{dr}{dt}\right) + \pi\!\left(R^2-r^2\right)\frac{dH}{dt} = 0.
\label{eq:const_volume_constraint}
\end{equation}

\paragraph{Analytical servo.}
The four conditions above---radial (\cref{eq:cond1_radial}), axial (\cref{eq:cond2_axial}), torsional (\cref{eq:cond3_torsion_error}), and constant-volume (\cref{eq:const_volume_constraint})---provide four independent equations for the four kinematic unknowns ($dr/dt$, $dR/dt$, $dH/dt$, $d\theta/dt$). The system is therefore closed and can be solved analytically.

An analytical solution is obtained by reducing the four unknowns to two through two successive eliminations:
\begin{enumerate}
\item \emph{Torsion.} The torsional condition (\cref{eq:cond3_torsion_error}) does not enter the kinematic constraints that govern $dr/dt$, $dR/dt$, and $dH/dt$. It is therefore solved independently at each timestep by adjusting the rotation angle rate:
\begin{equation}
\frac{d\theta}{dt} = \frac{e_T}{\Delta t}\,S_T,
\label{eq:cond3_torsion}
\end{equation}
where $\Delta t$ is the simulation timestep and $S_T$ is the torsional servo coefficient, determined from the contact-stiffness-weighted rotational stiffness $I_{rot}$ about the specimen axis (Fig.~\ref{fig:servo_geometry}):
\begin{equation}
I_{rot} = \sum k_n\bigl[(\mathbf{r}_d\times\hat{\mathbf{z}})\cdot\mathbf{n}\bigr]^2, \qquad S_T = \frac{1}{I_{rot}},
\label{eq:irot_scs}
\end{equation}
where $\mathbf{r}_d$ is the position vector from the rotation axis to the contact point in the horizontal plane, $\hat{\mathbf{z}}$ is the unit vector along the specimen axis, $\mathbf{n}$ is the contact normal directed from the particle centre to the blade, and $k_n$ is the normal contact stiffness.

\begin{figure}
  \centering
  \begin{tikzpicture}[
    every node/.style={font=\small},
    lbl/.style={font=\footnotesize},
    dimarr/.style={{Stealth[length=3pt]}-{Stealth[length=3pt]}, thin, black},
  ]
    \node[anchor=south west, inner sep=0] (img) at (0,0)
      {\includegraphics[width=0.5\linewidth]{servo_torsional_contact_geometry.png}};
    \begin{scope}[
      x={(img.south east)},
      y={(img.north west)},
    ]
      \coordinate (CP) at (0.646, 0.276);
      \coordinate (RA) at (0.4, 0.80);
      \coordinate (RA_h) at (0.4, 0.56);

      \node[lbl, red!70!black, above right] at (0.46, 0.4) {$\mathbf{n}$};

      \draw[-{Stealth[length=3pt]}, thick, blue!70!black] (RA_h) -- (CP);
      \node[lbl, blue!70!black, above] at ($(RA_h)!0.52!(CP)$) {$\mathbf{r}_d$};
      \draw[-{Stealth[length=3pt]}, thick, black] (RA_h) -- (RA);
      \node[lbl, black, right] at ($(RA_h)!0.7!(RA)$) {$\hat{\mathbf{z}}$};
    \end{scope}
  \end{tikzpicture}
  \caption{Contact geometry used to evaluate the rotational stiffness contribution for torsional servo control. $\mathbf{r}_d$: position vector from the rotation axis to the contact point in the horizontal plane; $\mathbf{n}$: contact normal directed from particle centre to blade; $\hat{\mathbf{z}}$: unit vector along the specimen axis.}
  \label{fig:servo_geometry}
\end{figure}
\item \emph{Volume constraint.} \Cref{eq:const_volume_constraint} expresses $dR/dt$ as a function of $dr/dt$ and $dH/dt$, reducing the system to two unknowns ($dr/dt$, $dH/dt$) with two target conditions ($e_r=0$ and $e_z=0$ in \cref{eq:cond1_radial,eq:cond2_axial}).
\end{enumerate}

After these two eliminations, both the radial and axial conditions depend on the same two unknowns ($dr/dt$, $dH/dt$). Because the volume constraint expresses $dR/dt$ as a function of $dr/dt$ and $dH/dt$, the outer-cylinder pressure $p_o^{\prime}$ is no longer independent: the radial condition, which involves both $p_i^{\prime}$ and $p_o^{\prime}$, is therefore affected by both $dr/dt$ and $dH/dt$ simultaneously. Similarly, the axial condition involves $p_z^{\prime}$ (affected by $dH/dt$) and $\sigma_r^{\prime}=({p_o^{\prime}R+p_i^{\prime}r})/({R+r})$ (Table~\ref{tab:hca_equations}), which depends on both cylinder movements. The two conditions are thus fully coupled. The next step is to quantify the contributions of $dr/dt$ and $dH/dt$ to the stress increments on all three walls, from which the coupled servo coefficients can be derived.

Let $K_{r,i}$ and $K_{r,o}$ denote the aggregate normal contact stiffness at the inner and outer cylinders, $A_i=2\pi rH$ and $A_o=2\pi RH$ the corresponding wall areas, and $K_z=\frac{1}{2}\sum_{c\in\mathrm{top/bottom}} k_n$ the effective axial contact stiffness with area $A_z=2\pi(R^2-r^2)$. Substituting \cref{eq:const_volume_constraint} into the pressure increments yields the coupled stiffness relation (derived in Appendix~\ref{sec:appendix_servo}):
\begin{equation}
\begin{pmatrix} \delta p_{dif,r}^{\prime} \\[4pt] \delta p_{dif,z}^{\prime} \end{pmatrix}
=
\underbrace{\begin{pmatrix} \hat{K}_{11} & \hat{K}_{12} \\[6pt] \hat{K}_{21} & \hat{K}_{22} \end{pmatrix}}_{\hat{\mathbf{K}}}
\begin{pmatrix} dr \\[4pt] dH \end{pmatrix},
\label{eq:coupled_stiffness}
\end{equation}
where
\begin{alignat}{2}
\hat{K}_{11} &= -\frac{K_{r,i}}{A_i}-\frac{r}{R}\,\frac{K_{r,o}}{A_o}, &\qquad
\hat{K}_{12} &= \frac{R^2-r^2}{2RH}\,\frac{K_{r,o}}{A_o}, \label{eq:Khat_main_row1}\\[4pt]
\hat{K}_{21} &= -\frac{K_{r,i}-\frac{r}{R}\,K_{r,o}}{A_i+A_o}, &\qquad
\hat{K}_{22} &= -\frac{K_z}{A_z}-\frac{(R^2-r^2)\,K_{r,o}}{2RH\,(A_i+A_o)}. \label{eq:Khat_main_row2}
\end{alignat}
The entries carry dimensions of Pa/m and incorporate the volume-constraint coupling. The determinant of $\hat{\mathbf{K}}$ is
\begin{equation}
\det(\hat{\mathbf{K}}) = \hat{K}_{11}\,\hat{K}_{22} - \hat{K}_{12}\,\hat{K}_{21}.
\label{eq:det_K_main}
\end{equation}
Setting $(\delta p_{dif,r}^{\prime},\,\delta p_{dif,z}^{\prime})=(-e_r,\,-e_z)$ and inverting \cref{eq:coupled_stiffness} yields the coupled servo equations:
\begin{align}
\frac{dr}{dt} &= S_{rr}\,\frac{e_r}{\Delta t} + S_{rz}\,\frac{e_z}{\Delta t}, \label{eq:servo_dr}\\[6pt]
\frac{dH}{dt} &= S_{zr}\,\frac{e_r}{\Delta t} + S_{zz}\,\frac{e_z}{\Delta t}, \label{eq:servo_dH}
\end{align}
where the servo coefficients are the entries of $\hat{\mathbf{K}}^{-1}$:
\begin{alignat}{2}
S_{rr} &= \frac{-\hat{K}_{22}}{\det(\hat{\mathbf{K}})}, &\qquad
S_{rz} &= \frac{\hat{K}_{12}}{\det(\hat{\mathbf{K}})}, \label{eq:servo_coeff_r}\\[4pt]
S_{zr} &= \frac{\hat{K}_{21}}{\det(\hat{\mathbf{K}})}, &\qquad
S_{zz} &= \frac{-\hat{K}_{11}}{\det(\hat{\mathbf{K}})}. \label{eq:servo_coeff_z}
\end{alignat}
The diagonal terms $S_{rr}$ and $S_{zz}$ represent the direct servo response, while the off-diagonal terms $S_{rz}$ and $S_{zr}$ capture the cross-coupling between the radial and axial conditions. The remaining kinematic variable $dR/dt$ is then obtained by substituting $dr/dt$ and $dH/dt$ into \cref{eq:const_volume_constraint}. The exact analytical solution of the coupled stiffness matrix ensures that the servo coefficients are physically consistent, providing the correct basis on which a relaxation factor ($<1$) can be applied for numerical stability. Under the constant specimen height condition ($dH/dt=0$) adopted by \citet{Han2024jgs}, only the radial condition remains active and the servo coefficient reduces to
\begin{equation}
S_{rr} = -\frac{1}{\hat{K}_{11}} = \frac{1}{\dfrac{K_{r,i}}{A_i}+\dfrac{r}{R}\,\dfrac{K_{r,o}}{A_o}}.
\label{eq:Srr_constantH_main}
\end{equation}

In undrained loading, the contact network thins as the effective confining stress decays, and the aggregate wall stiffnesses $K_{r,i}$, $K_{r,o}$ and the end-platen stiffness $K_z=\tfrac{1}{2}(K_z^{top}+K_z^{bot})$ that enter $\hat{\mathbf{K}}$ can drop towards zero, driving $\det(\hat{\mathbf{K}})$ towards zero and the servo coefficients towards arbitrarily large values. To prevent near-singular behaviour in the closed-form inversion, each per-wall aggregate contact stiffness is floored at $10\%$ of its end-of-consolidation value before $\hat{\mathbf{K}}$ is assembled, which keeps every entry of $\hat{\mathbf{K}}$ bounded away from zero by construction. An analogous $20\%$ floor is applied to the rotational stiffness $I_{rot}$ in the independent torsional servo (\cref{eq:irot_scs}), which is more sensitive to the loss of blade--particle contacts as $\sigma_z^{\prime}\!\to\!0$. Both floors function as safety margins on the closed-form inversion and are exercised explicitly in Section~\ref{subsec:numerical_verification}.

\begin{figure}
  \centering
  \begin{tikzpicture}[
    node distance=10pt,
    box/.style={draw, rounded corners=2pt, minimum width=0.88\linewidth,
                minimum height=7mm, align=center, font=\small},
    decision/.style={draw, diamond, aspect=2.5, minimum width=0.7\linewidth,
                     align=center, font=\small, inner sep=1pt},
    arr/.style={-{Stealth[length=5pt]}, thick},
    label/.style={font=\footnotesize\itshape, text=black!70},
  ]
    \node[box] (step1)
      {Compute $p_i^{\prime},\;p_o^{\prime},\;p_z^{\prime},\;T$
       from particle--wall contacts};
    \node[box, below=of step1] (step2)
      {Evaluate errors:\;
       $e_r=p_{dif,r}^{\prime}-p_{dif,r}^{tar}$,\;
       $e_z=p_{dif,z}^{\prime}-p_{dif,z}^{tar}$,\;
       $e_T=T-T^{tar}$};
    \node[box, below=of step2] (step3)
      {Torsion (independent):\;
       $d\theta/dt = (e_T/\Delta t)\,S_T$};
    \node[box, below=of step3] (step4)
      {Assemble $\hat{\mathbf{K}}$ from contact stiffnesses;\;
       compute $S_{rr},\,S_{rz},\,S_{zr},\,S_{zz}$};
    \node[box, below=of step4] (step5)
      {Coupled radial--axial servo:\;
       $dr/dt$,\; $dH/dt$ \;(\cref{eq:servo_dr,eq:servo_dH})};
    \node[box, below=of step5] (step6)
      {Volume constraint $\rightarrow$ $dR/dt$
       \;(\cref{eq:const_volume_constraint})};
    \node[box, below=of step6] (step7)
      {Move walls;\; advance timestep};

    \draw[arr] (step1) -- (step2);
    \draw[arr] (step2) -- (step3);
    \draw[arr] (step3) -- (step4);
    \draw[arr] (step4) -- (step5);
    \draw[arr] (step5) -- (step6);
    \draw[arr] (step6) -- (step7);
    \draw[arr] (step7.south) -- ++(0,-0.3) -| ([xshift=8pt]step1.east |- step7.south) -- ([xshift=8pt]step1.east) -- (step1.east)
      node[label, pos=0.45, right=20pt, rotate=90, anchor=south] {next timestep};
  \end{tikzpicture}
  \caption{Servo procedure at each simulation timestep during undrained loading.}
  \label{fig:servo_flowchart}
\end{figure}

The excess pore water pressure is inferred from the difference between the mean radial total and effective stresses:
\begin{equation}
u = \sigma_r - \sigma_r^{\prime},
\label{eq:epwp_definition}
\end{equation}
where $\sigma_r$ is the mean radial total stress prescribed by the laboratory boundary conditions (determined by $p_o$ and $p_i$) and $\sigma_r^{\prime}$ is the current mean radial effective stress evaluated from the DEM contact forces (Table~\ref{tab:hca_equations}). For the constant-cell-pressure cases studied here, $\sigma_r$ is time-invariant and equal to $\sigma_{r0}^{\prime}$, so \cref{eq:epwp_definition} reduces to $u = \sigma_{r0}^{\prime} - \sigma_r^{\prime}$. The ratio $\ru$ is defined as $\ru = u/\sigma_{r0}^{\prime}$. Figure~\ref{fig:servo_flowchart} summarizes the complete servo procedure applied at each simulation timestep.

Compared with earlier DEM-based HCA servo strategies, the present formulation makes three contributions explicit. First, the governing equations are systematically reduced through elimination of the volume constraint and torsional independence before the coupled system is solved, clarifying the mathematical structure of the problem. Second, the laboratory-to-DEM mapping is established through two pressure differences ($p_{dif,r}^{\prime}$ and $p_{dif,z}^{\prime}$), treated as general time-dependent targets $p_{dif,r}^{tar}(t)$ and $p_{dif,z}^{tar}(t)$ that are invariant to the unknown pore water pressure and reproduce the independent $p_o$, $p_i$, $p_z$ control available on a laboratory HCA. Third, the cross-coupling between the radial and axial degrees of freedom is quantified analytically through the off-diagonal entries of $\hat{\mathbf{K}}$, so that the servo coefficients correctly reflect the mutual influence of $dr/dt$ and $dH/dt$ on both boundary conditions. None of these steps were formalized in earlier work: \citet{Han2024jgs} solved only the radial and torsional conditions while keeping the specimen height fixed ($dH/dt=0$) and the radial pressure-difference target held at $p_{dif,r}^{tar}=0$ (i.e., $p_o=p_i$), and \citet{Ma2024} introduced an axial servo but derived each servo coefficient from the contact stiffness at the corresponding wall independently, without reflecting the cross-coupling among the wall movements.

\subsection{Parameter calibration and validation}
\label{subsec:calibration_validation}

The micro-parameters listed in Table~\ref{tab:dem_parameters} were calibrated against the undrained monotonic torsional shear data of \citet{Nakata1998} for very dense Toyoura sand ($D_r \approx 90\%$, $\pmean_0=\SI{100}{kPa}$, \Kzero{} = 1.0) \citep[see also][]{Ma2024}. Because the micro-parameters collectively influence shear stiffness, dilatancy, and liquefaction resistance, the calibration targeted these behaviors in sequence. Starting from a base parameter set at a given void ratio, the contact Young's modulus $E$ and rolling friction coefficient $\mu_r$ were adjusted while keeping the normal-to-shear stiffness ratio $\kappa$ constant, aiming to reproduce the shear stiffness observed in the $\varepsilon_q$--$q$ response of the monotonic shear tests. The effective stress path in $\pmean$--$q$ space was then examined to assess dilatancy. If the simulated dilatancy was insufficient---i.e., the mean effective stress at phase transformation was lower than the experimental value---the specimen was further compacted to a lower void ratio, which increased both dilatancy and shear stiffness; $E$ was then reduced accordingly to restore the target shear stiffness. This iterative process between void ratio, contact modulus, and rolling friction was repeated until both the shear stiffness and dilatancy simultaneously matched the monotonic experimental data. During the monotonic calibration, the axial pressure difference was maintained at zero ($p_{dif,z}^{tar}=0$ in \cref{eq:cond2_axial}) to ensure $\sigma_z^{\prime}=\sigma_r^{\prime}$ during shearing, and shear was imposed through prescribed top-platen rotation at $\omega = \SI{2}{rad/s}$, well within the quasi-static regime quantified in Section~\ref{subsec:numerical_verification}.

The calibrated parameter set was then validated against the cyclic liquefaction resistance data of \citet{Ishihara1985} for the same sand under identical conditions. Specimens were subjected to sinusoidal torsional shear at five CSR levels (0.20, 0.25, 0.30, 0.35, and 0.40) under constant specimen height ($dH/dt=0$), matching the experimental protocol of \citet{Ishihara1985}, with liquefaction defined by single-amplitude shear strain $\gamma_{SA}=2.5\%$.

\begin{figure*}
  \centering
  \safeincludegraphics[width=\textwidth]{validate_combined.png}
  \caption{Calibration and validation of the DEM-HCA model for very dense Toyoura sand ($D_r \approx 90\%$, $\pmean_0=\SI{100}{kPa}$, \Kzero{} = 1.0): (a) effective stress path and (b) stress--strain relationship from undrained monotonic shear, calibrated against \citet{Nakata1998}; (c) cyclic stress ratio versus number of cycles to liquefaction, validated against \citet{Ishihara1985}.}
  \label{fig:validation_combined}
\end{figure*}

Fig.~\ref{fig:validation_combined} summarizes the calibration and validation results. The effective stress path (Fig.~\ref{fig:validation_combined}a) captures the initial contraction followed by dilation towards the critical state line, matching the experimental trend. The stress--strain relationship (Fig.~\ref{fig:validation_combined}b) reproduces the shear stiffness and strain-hardening behavior. These two aspects confirm a successful calibration against the monotonic data. The liquefaction resistance curve (Fig.~\ref{fig:validation_combined}c) shows good agreement with the cyclic data across the full range of CSR, validating the calibrated parameters under cyclic loading conditions without further adjustment. The simultaneous agreement across stiffness, dilatancy, and cyclic liquefaction resistance confirms the reliability of the micro-parameters (Table~\ref{tab:dem_parameters}) and demonstrates that the servo mechanism performs well under both axial stress-control (monotonic calibration) and displacement-control (cyclic validation) modes.

\paragraph{Critical-state determination from large-strain monotonic undrained shear.}
The critical-state line of the calibrated specimen is determined from a separate undrained monotonic torsional-shear simulation on the same $D_r \approx 90\%$ specimen at $\pmean_0 = \SI{100}{kPa}$ and \Kzero{} = 1.0, continued into the regime where $q/p^{\prime}$ stops evolving with strain. The resulting effective stress path and stress-ratio evolution are shown in Fig.~\ref{fig:csl_determination}. The stress ratio $q/p^{\prime}$ rises to a dilatant peak of $\sim 1.27$ near $\varepsilon_q \approx 6\%$, then descends through the phase-transformation transition and locks onto a clean plateau over $\varepsilon_q \in [20, 40]\%$, from which the critical-state slope is fit at $M_{cs} = 0.874 \pm 0.018$, corresponding to a friction angle $\varphi_{cs} = 30.3^{\circ}$ via the Mohr--Coulomb relation $M_{cs} = \sqrt{3}\sin\varphi_{cs}$ (which holds here because the HCA torsional boundary conditions impose pure-shear loading by construction, intermediate principal stress ratio $b = 0.5$). Comparison with the laboratory literature is most cleanly made through the friction angle, since the canonical drained-triaxial-compression dataset of \citet{VerdugoIshihara1996} for Toyoura sand is at $b = 0$ rather than $b = 0.5$: they report a steady-state slope of $0.62$ in the $q/2$--$p^{\prime}$ plane, equivalent to $M_{TC} = q/p^{\prime} = 1.24$ and a friction angle $\varphi_{ss} \approx 31^{\circ}$. Our $\varphi_{cs} = 30.3^{\circ}$ agrees with this within $0.7^{\circ}$, and the triaxial-compression slope predicted from our $\varphi_{cs}$, $M_{TC} = 6\sin\varphi_{cs}/(3 - \sin\varphi_{cs}) = 1.21$, matches their measured $M_{TC} = 1.24$ within $\sim 2\%$. The $M_{cs} = 0.874$ reported here is therefore the monotonic, pure-shear ($b = 0.5$) critical-state slope; real granular materials exhibit a weak dependence of $\varphi_{cs}$ on the intermediate principal stress ratio that lies beyond the scope of this calibration.

\begin{figure*}
  \centering
  \safeincludegraphics[width=0.733\textwidth]{critical_state.png}
  \caption{Critical-state determination on the calibrated $D_r \approx 90\%$ specimen ($\pmean_0 = \SI{100}{kPa}$, \Kzero{} = 1.0): (a)~effective stress path in $\pmean$--$q$ space, with the red segment marking the critical-state range $\varepsilon_q \in [20, 40]\%$ used to fit $M_{cs} = 0.874$; (b)~stress ratio $q/p^{\prime}$ versus deviatoric strain $\varepsilon_q$, settling onto the plateau at $q/p^{\prime} = 0.874 \pm 0.018$.}
  \label{fig:csl_determination}
\end{figure*}

\subsection{Numerical verification of the coupled servo}
\label{subsec:numerical_verification}

The validation in Section~\ref{subsec:calibration_validation} establishes that the calibrated parameters reproduce the macroscopic response of laboratory Toyoura sand. A separate question is whether the servo formulation itself solves the boundary conditions stated in \cref{eq:cond1_radial,eq:cond2_axial,eq:cond3_torsion_error,eq:const_volume_constraint} faithfully through the regimes of interest, and whether the cyclic loading is conducted slowly enough that grain-inertia effects do not contaminate the response. This subsection reports three verification quantities measured throughout a representative high-CSR cyclic run on the dense specimen ($D_r \approx 75\%$, $\Kzero{} = 1.0$, CSR $= 0.400$) under the more demanding constant-total-stress axial boundary ($p_{dif,z}^{tar}$ held constant, $dH/dt$ free), so that the full $2\times 2$ coupled servo is exercised.

\begin{figure*}
  \centering
  \safeincludegraphics[width=\textwidth]{numerical_verification.png}
  \caption{Numerical verification of the coupled servo on a representative high-CSR cyclic run ($D_r \approx 75\%$, $\Kzero{} = 1.0$, CSR $= 0.400$, constant-total-axial-stress boundary): (a)~specimen-volume drift $V/V_0 - 1$ confirming exact enforcement of the constant-volume constraint (\cref{eq:const_volume_constraint}); (b)~tracking errors of the radial and axial pressure-difference targets (\cref{eq:cond1_radial,eq:cond2_axial}); (c)~inertial number $I$ on log-scale; horizontal references at $I=10^{-3}$ and $I=10^{-2}$ are the strict and representative quasi-static thresholds (see text), and the vertical dashed line marks the onset of liquefaction.}
  \label{fig:numerical_verification}
\end{figure*}

The constant-volume constraint is enforced kinematically through the analytical relation between $dr/dt$, $dR/dt$, and $dH/dt$ (\cref{eq:const_volume_constraint}), so the specimen volume is expected to remain exactly conserved up to floating-point round-off. Fig.~\ref{fig:numerical_verification}a confirms this: the relative volume drift $V(t)/V_0 - 1$ stays below $10^{-7}$ throughout the run, which is several orders of magnitude smaller than the laboratory drift attributable to membrane compliance and water compressibility. Fig.~\ref{fig:numerical_verification}b shows the corresponding tracking errors for the radial and axial pressure-difference targets ($e_r = p_{dif,r}^{\prime} - p_{dif,r}^{tar}$, $e_z = p_{dif,z}^{\prime} - p_{dif,z}^{tar}$): both grow modestly as the effective confining stress decays towards liquefaction, but remain bounded throughout, with peak magnitudes a few percent of the initial mean effective stress. The closed-form $2 \times 2$ inversion of $\hat{\mathbf{K}}$ therefore continues to deliver well-conditioned wall velocities even when the underlying contact network is severely depleted; in particular, the $10\%$ stiffness floor and $20\%$ moment-of-inertia floor introduced in Section~\ref{subsec:cyclic_loading} are sufficient to keep $\det\hat{\mathbf{K}}$ bounded away from zero throughout the post-liquefaction excursions.

The inertial number $I=|\dot\gamma_{z\theta}|\,d_{50}\sqrt{\rho_s/\pmean}$, with $d_{50}=\SI{2.0}{mm}$ and $\rho_s=\SI{2650}{kg/m^3}$, plotted in Fig.~\ref{fig:numerical_verification}c quantifies the strength of grain-inertia effects relative to the prevailing confining stress. \citet{GDRMiDi2004} report that the effective friction becomes shear-rate independent below $I = 10^{-3}$ in steady plane shear, defining a strict quasi-static threshold; \citet{daCruzEmam2005} subsequently use $I = 10^{-2}$ as the representative quasi-static value. In the present run, the representative value $I < 10^{-2}$ is satisfied for the entire pre-liquefaction stage ($N_c/N_L < 1$, $\ru < 0.95$), and the strict threshold $I < 10^{-3}$ is satisfied throughout $96.9\%$ of that stage. The remaining $\sim 3\%$ excursion above $10^{-3}$ is confined to the final cycle before liquefaction, where the mean effective stress has already decayed by more than an order of magnitude and the loading enters the rapid-transition stage of \citet{Wang2025inertial}. After liquefaction, $I$ rises further to peak values of order $10^{-2}$, an unavoidable physical consequence of $\pmean \to 0$ in the inertial-number definition that is shared with laboratory undrained cyclic tests \citep{Wang2025inertial}; the strict quasi-static interpretation of $I$ does not apply in this regime in either the simulation or the laboratory. The monotonic-shear runs of Section~\ref{subsec:calibration_validation} (at the loading rate stated there) likewise satisfy both quasi-static thresholds, since their stress paths never approach the $\pmean \to 0$ regime that drives $I$ upward in the present cyclic case.

\subsection{Microscopic descriptors}
\label{subsec:micro_descriptors}

While the macroscopic response characterizes the overall liquefaction behavior, a key advantage of DEM is its ability to access particle-scale information that is not directly measurable in laboratory tests. To exploit this advantage, several microscopic descriptors are introduced to probe the contact network from complementary perspectives: the mechanical coordination number \Zm{}, the count- and force-weighted fabric tensors $\Phi_{ij}$ and $\Phi^f_{ij}$ together with their scalar invariants $(A, B)$, the signed anisotropy indicator $\alpha$, and the contact density $\rho_c$.

The mechanical coordination number \citep{Thornton2000,Oda1977} is defined as
\begin{equation}
Z_m = \frac{2N_c - N_1}{N_p - N_1 - N_0},
\label{eq:zm_definition_method}
\end{equation}
where $N_c$ is the total number of interparticle contacts, $N_p$ is the total number of particles, and $N_1$ and $N_0$ are the numbers of particles with one and zero contacts, respectively. Floaters are excluded to focus on the load-bearing contact network. $Z_m$ serves as a scalar measure of packing compactness.

The directionality of the contact network is characterized through the fabric tensor \citep{Oda1982fabric,Oda1982stat}:
\begin{equation}
\Phi_{ij} = \frac{1}{N_c}\sum_c n_i^c\,n_j^c,
\label{eq:fabric_tensor}
\end{equation}
where $n_i^c$ and $n_j^c$ are the components of the unit contact normal at contact $c$. Its force-weighted counterpart \citep{Zhou2022,Song2024},
\begin{equation}
\Phi^f_{ij} = \frac{\sum_c f_n^c\, n_i^c\, n_j^c}{\sum_c f_n^c},
\label{eq:fabric_tensor_force}
\end{equation}
weights each contact by its normal-force magnitude $f_n^c$ and therefore reflects the directional partitioning of the load-bearing subnetwork rather than the contact topology alone; the two diagonals coincide only when contact forces are uniform across orientations. Because the HCA specimen is axisymmetric about the vertical axis $z$, both tensors are expressed in the cylindrical coordinate system $(r,\theta,z)$ in what follows. From each tensor we extract two scalar invariants: the deviator magnitude
\begin{equation}
A = \sqrt{\tfrac{3}{2}}\,\bigl\|\Phi - I/3\bigr\|_F \in [0,1],
\label{eq:fabric_invariants}
\end{equation}
and the tilt $B = \arccos|\hat{\bm v}_1 \cdot \hat{z}|$ of the largest eigenvector $\hat{\bm v}_1$ from the axial direction; subscripts $c$ and $n$ identify the count- and normal-force-weighted versions, so that $(A_c, B_c)$ are computed from $\Phi$ and $(A_n, B_n)$ from $\Phi^f$. The signed scalar indicator $\alpha$ \citep{YimsiriSoga2010,Otsubo2022} is defined from the axial diagonal component $\Phi_{zz}$:
\begin{equation}
\alpha = \frac{5(3\Phi_{zz} - 1)}{5\Phi_{zz} + 1}.
\label{eq:alpha_definition_method}
\end{equation}
Positive $\alpha$ indicates that contact normals concentrate toward the vertical axis (compression-state anisotropy, \Kzero{} $< 1.0$), whereas negative $\alpha$ indicates concentration toward the horizontal plane (extension-state anisotropy, \Kzero{} $> 1.0$).

The contact density \citep{Han2023} describes the average number of contacts per unit solid angle on the unit sphere of contact normal orientations:
\begin{equation}
\rho_c(\theta_z,\,\phi_{\mathrm{cir}}) = \frac{2\,N(\theta_z,\,\phi_{\mathrm{cir}})}{N_p\,\Delta\Omega},
\label{eq:contact_density}
\end{equation}
where $\theta_z$ and $\phi_{\mathrm{cir}}$ are the polar and azimuthal angles in the spherical coordinate system, $N(\theta_z,\,\phi_{\mathrm{cir}})$ is the number of contacts whose normals fall within the angular bin of widths $\Delta\theta_z$ and $\Delta\phi_{\mathrm{cir}}$, and $\Delta\Omega = (\cos\phi_{\mathrm{cir}} - \cos(\phi_{\mathrm{cir}}+\Delta\phi_{\mathrm{cir}}))\,\Delta\theta_z$ is the solid angle subtended by that bin. Unlike the fabric tensor, $\rho_c$ retains the full directional information and thus provides a more complete picture of the evolving contact network.

Together, these descriptors probe the contact network across a hierarchy of directional detail---from the scalar packing measure \Zm{} through tensorial anisotropy ($\Phi$, $\Phi^f$ and their invariants) to the full directional field $\rho_c$---and connect particle-scale fabric and force evolution to the macroscopic liquefaction resistance trends reported in \cref{sec:results}.
